Im Youtube-Video \cite{buch:feynman} wird vorgeschlagen, das Integral
\[
\int \frac{x^2}{(1+x^2)^2}\,dx
\]
mit der Feynman-Methode zu lösen, die darauf beruht, einen Parameter
$a$ in das Integral
\[
I(a)
=
\int \frac{1}{1+ax^2}\,dx
\]
einzuführen und dann nach $a$ abzuleiten.
Durch Integration unter dem Integralzeichen wird daraus
\[
I'(a)
=
-
\int \frac{x^2}{(1+ax^2)^2}\,dx,
\]
d.~h.~das gesuchte Integral ist $I'(a)$.
Da
\[
I(a)
=
\int \frac{1}{1+ax^2}\,dx
=
\int \frac{1}{1+(\!\sqrt{a}x)^2}\,\!\sqrt{a}\,dx
=
\!\sqrt{a}
\arctan(\!\sqrt{a}x)
\]
ein wohlbekanntes Integral ist, muss man nur noch die Ableitung ausrechnen
und findet
\[
-
I'(1) = \frac12 \arctan(x) - \frac{x}{2(1+x^2)}.
\]
Bestimmen Sie das Integral stattdessen mit dem in diesem Kapitel
beschriebenen Algorithmus.

\begin{loesung}
Die quadratfreie Faktorisierung des Nenners ist $q(x)^2$ mit $q(x)=1+x^2$.
Wir müssen also als erstes die Methode von Hermite anzuwenden, um den
Exponenten im Nenner zu reduzieren.
Dazu müssen $s(x)$ und $t(x)$ gefunden werden, so dass
\[
\begin{array}{rcrcl}
s(x)q(x)    &+& t(x)q'(x) &=& x^2 \\
s(x)(1+x^2) &+& t(x) 2x &=& x^2
\end{array}
\]
gilt.
Mit dem euklidischen Algorithmus findet man
\[
\underbrace{x^2}_{\displaystyle=s(x)}\mathstrut\cdot q(x)
\underbrace{\mathstrut-\frac{x^3}{2}}_{\displaystyle=t(x)}\mathstrut\cdot q'(x) = x^2
\]
und damit
\begin{align*}
\int
\frac{x^2}{(1+x^2)^2}\,dx
&=
\int \frac{x^2}{q(x)}\,dx
-
\int \frac{x^3}{2} \frac{q'(x}{q(x)^2}\,dx
\\
&=
\int \frac{x^2}{q(x)}\,dx
-
\frac{x^3}{2} \frac{-1}{q(x)}
-
\int
\frac{3x^2}{2} \frac{-1}{q(x)}
\,dx
\\
&=
\frac{x^3}{2(1+x^2)}
-
\frac{1}{2}
\int
\frac{x^2}{q(x)}\,dx
\\
&=
\frac{x^3}{2(1+x^2)}
-
\frac{1}{2}
\int
1-
\frac{1}{q(x)}
\,dx
\\
&=
\frac{x^3}{2(1+x^2)}
-
\frac{x}{2}
+
\frac12
\int\frac{1}{1+x^2}\,dx
\\
&=
-
\frac12
\frac{x}{1+x^2}
+
\frac12\arctan x.
\qedhere
\end{align*}
\end{loesung}
